\documentclass{SibFU-docs}

% доп. пакеты
\usepackage{graphicx}
\usepackage{float}
\usepackage{hyperref}
\usepackage{caption}
\usepackage{subcaption}
\usepackage{amsmath}
% предотвартим ошибку с \Bbbk перед загрузкой пакета amssymb
\let\Bbbk\relax
\usepackage{amssymb}
\usepackage{booktabs}
\usepackage{tabularx}
\usepackage{longtable}
\usepackage{enumitem}
\usepackage{multirow}

% гиперссылки
\hypersetup{
    colorlinks=true,
    linkcolor=black,
    citecolor=blue,
    urlcolor=blue
}


% библиография
\addbibresource{report.bib}

% титульный лист
\begin{document}
\setlist[itemize]{left=1.3cm, itemsep=1pt, topsep=2pt} % для выравнивания списков
\setlist[enumerate]{left=1.3cm, itemsep=1pt, topsep=2pt}
\renewcommand{\contentsname}{}
\mycovertitle
{Гуманитарный институт}
{Кафедра прикладной информатики в искусстве и интерактивном медиа}
{Табличные данные и базы данных (реляционные и нереляционные, СУБД)}
{ст. преподаватель И.~Р.~Нигматуллин}
{Захаров И.~М., Логинова Ю.~В., Тетерина П.~А.}
 
\begin{center}
\bfseries РЕФЕРАТ
\end{center}
\noindent

Реферат по теме: «Табличные данные и базы данных (реляционные и нереляционные, СУБД)» содержит 37 страниц и 18 использованных источников.

Ключевые слова: БАЗЫ ДАННЫХ, СИСТЕМЫ УПРАВЛЕНИЯ БАЗАМИ ДАННЫХ (СУБД), РЕЛЯЦИОННАЯ МОДЕЛЬ, НЕРЕЛЯЦИОННЫЕ БАЗЫ ДАННЫХ (NoSQL), ЯЗЫК SQL, НОРМАЛИЗАЦИЯ, ЦЕЛОСТНОСТЬ ДАННЫХ, ОБЛАЧНЫЕ БАЗЫ ДАННЫХ, ТРАНЗАКЦИИ, МАСШТАБИРУЕМОСТЬ, БЕЗОПАСНОСТЬ ДАННЫХ

В работе проводится комплексное исследование современных подходов к организации, хранению и управлению структурированной информацией в базах данных. Рассматриваются как классические реляционные системы, так и современные нереляционные модели, их архитектурные особенности, области применения и сравнительные характеристики.

Детально проанализированы фундаментальные принципы организации данных, включая процесс нормализации, обеспечивающий устранение избыточности и поддержание целостности информации. Особое внимание уделено системам управления базами данных (СУБД), их ключевым функциям: управлению транзакциями с гарантиями ACID, обеспечению безопасности, оптимизации производительности и механизмам восстановления после сбоев.

Отдельный раздел посвящен языку SQL как универсальному инструменту взаимодействия с реляционными базами данных, его возможностям для сложных запросов и аналитических операций. Рассмотрены современные концепции надежности и целостности данных, включая теорему CAP, модели согласованности в распределенных системах и стратегии обеспечения отказоустойчивости.

Завершающая часть работы охватывает облачные подходы к работе с базами данных, включая модель «база данных как сервис» (DBaaS), специализированные облачные предложения для различных типов СУБД и онлайн-инструменты для разработки и администрирования.

Основные выводы:
\begin{enumerate}
\item Современные информационные системы требуют осознанного выбора между реляционными и нереляционными моделями данных на основе специфики задач: реляционные системы обеспечивают строгую целостность и сложные запросы, а нереляционные предлагают горизонтальную масштабируемость и гибкость для специализированных сценариев;
\item Системы управления базами данных эволюционировали в сложные программные комплексы, предоставляющие не только хранение информации, но и комплексные механизмы обеспечения транзакционности, безопасности, производительности и отказоустойчивости, что делает их фундаментом для критически важных приложений;
\item Язык SQL остается универсальным стандартом для работы со структурированными данными, постоянно развиваясь и интегрируясь с современными технологиями больших данных, облачными платформами и аналитическими системами;
\item Переход к облачным моделям работы с базами данных (DBaaS) кардинально меняет парадигму управления информацией, предлагая экономическую эффективность, автоматическое масштабирование и снижение операционной нагрузки, что особенно актуально в условиях цифровой трансформации бизнеса и научных исследований.
\end{enumerate}

Рекомендации:
\begin{itemize}[label=-]
\item При проектировании новых информационных систем проводить тщательный анализ требований к данным для обоснованного выбора между реляционными и нереляционными моделями, учитывая не только текущие потребности, но и стратегию развития системы;
\item Осваивать и применять современные возможности SQL, включая оконные функции, рекурсивные запросы и интеграцию с полуструктурированными данными, для повышения эффективности работы с информацией в реляционных системах;
\item Внедрять облачные базы данных и DBaaS-решения для новых проектов и постепенно мигрировать существующие системы для снижения затрат на инфраструктуру и администрирование, уделяя особое внимание вопросам безопасности и соответствия регуляторным требованиям;
\item Разрабатывать и регулярно тестировать комплексные стратегии обеспечения целостности, надежности и восстановления данных, включая определение метрик RPO/RTO, реализацию репликации, резервного копирования и процедур аварийного восстановления.
\end{itemize}

\newpage

\thispagestyle{empty}
\begin{center}
\bfseries СОДЕРЖАНИЕ
\end{center}
\vspace{-6pt}
\tableofcontents
\clearpage

\section*{}
\vspace*{-2.5em}
\begin{center}\textbf{ВВЕДЕНИЕ}\end{center}
\phantomsection
\addcontentsline{toc}{section}{Введение}

Базы данных являются фундаментальным компонентом современной цифровой инфраструктуры, лежащим в основе практически всех информационных систем — от банковских приложений и электронной коммерции до научных исследований и государственного управления. Они обеспечивают структурированное хранение, эффективный доступ и целостное управление информацией, преобразуя необработанные данные в стратегический актив организаций. Однако кажущаяся распространенность и доступность технологий баз данных часто маскирует сложность их грамотного проектирования, выбора и эксплуатации, требующей глубокого понимания архитектурных принципов и компромиссов.

В условиях экспоненциального роста объемов информации и разнообразия форматов данных особенно остро встают вопросы выбора адекватной модели хранения, обеспечения производительности и масштабируемости систем. Противоречивые требования различных приложений — необходимость строгой целостности транзакций в финансовых системах, горизонтальная масштабируемость для веб-приложений с миллионами пользователей, гибкость для работы с быстро изменяющимися данными — приводят к сосуществованию различных парадигм: классических реляционных и современных нереляционных систем. Непонимание их принципиальных различий и областей эффективного применения ведет к выбору неоптимальных решений, проблемам с производительностью, сложностям масштабирования и, в конечном итоге, к снижению эффективности информационных систем.

Современное состояние вопроса характеризуется богатым набором технологий и подходов: от проверенных временем реляционных СУБД с их развитыми механизмами транзакционности и стандартизированным языком SQL до разнообразных NoSQL-систем, оптимизированных для конкретных сценариев использования. Развитие облачных вычислений привело к появлению модели «база данных как сервис», которая меняет парадигму управления информационной инфраструктурой. Однако эффективное использование этого многообразия требует системного понимания архитектурных принципов, ограничений и компромиссов каждой технологии.

Актуальность темы обусловлена всеобщей цифровой трансформацией, в рамках которой данные становятся ключевым активом организаций. Грамотное проектирование и управление системами хранения данных напрямую влияет на эффективность бизнес-процессов, качество принимаемых решений и конкурентные преимущества. Новизна работы заключается в комплексном рассмотрении всего спектра технологий баз данных — от фундаментальных принципов организации информации и различий между моделями данных до практических аспектов использования SQL, обеспечения целостности и надежности, а также современных облачных подходов.

Цель работы: исследовать совокупность принципов, моделей и технологий, составляющих современную экосистему баз данных, и проанализировать их применение в различных сценариях.

Задачи:
\begin{enumerate}
\item Изучить общие принципы организации и хранения данных в базах данных, включая процесс нормализации и методы проектирования;
\item Проанализировать различия между реляционными и нереляционными моделями данных, их архитектурные особенности, преимущества, недостатки и типичные области применения;
\item Рассмотреть основные функции и возможности систем управления базами данных, включая управление транзакциями, обеспечение безопасности, оптимизацию производительности и механизмы восстановления;
\item Исследовать язык SQL как инструмент взаимодействия с базами данных, его возможности для сложных запросов и интеграцию с современными технологиями;
\item Проанализировать концепции целостности и надежности в системах хранения данных, включая теорему CAP, модели согласованности и стратегии обеспечения отказоустойчивости;
\item Определить современные облачные подходы и онлайн-инструменты для работы с базами данных, их преимущества и перспективы развития.
\end{enumerate}

Методы: системный анализ архитектурных подходов, сравнительный анализ технологий и моделей данных, обобщение и структурирование информации из профессиональной литературы, стандартов и лучших практик проектирования и эксплуатации баз данных.

\newpage
% вручную увеличить номер секции и записать в оглавление с номером,
% чтобы в теле не печатался автоматический номер, но в tableofcontents он остался
\refstepcounter{section}
\addcontentsline{toc}{section}{\protect\numberline{\thesection} Общие принципы организации и хранения данных в базах данных}
\vspace*{-2.5em}
\begin{center}\textbf{\thesection\ ОБЩИЕ ПРИНЦИПЫ ОРГАНИЗАЦИИ И ХРАНЕНИЯ ДАННЫХ В БАЗАХ ДАННЫХ}\end{center}
\vspace*{-1.5em}

\ManualSubsection{Основы проектирования и хранения данных в базах}

Базы данных представляют собой не просто хранилища информации, а сложные системы, организованные по определенным принципам, которые обеспечивают эффективное управление данными на протяжении всего их жизненного цикла. Эти принципы формируют фундамент, на котором строятся все современные системы хранения и обработки информации, от простых настольных приложений до распределенных облачных платформ.

\ManualSubsubsection{Эволюция подходов к организации данных}

Исторически развитие методов организации данных прошло несколько ключевых этапов. На начальном этапе информация хранилась в плоских файлах, что приводило к значительной избыточности и сложностям при обновлении. Появление иерархических и сетевых моделей в 1960-х годах позволило установить четкие связи между элементами данных, но эти модели оставались сложными в использовании и обслуживании. Настоящий прорыв произошел в 1970 году, когда Эдгар Кодд предложил реляционную модель данных, основанную на математической теории множеств и логике первого порядка \cite{codd1970}.

Реляционная модель принципиально изменила подход к организации данных, представив информацию в виде набора взаимосвязанных таблиц. Каждая таблица соответствует некоторой сущности предметной области и состоит из строк (записей) и столбцов (атрибутов). Строка представляет отдельный экземпляр сущности, а столбец — определенное свойство этой сущности. Связи между таблицами устанавливаются через ключевые поля, что позволяет избежать избыточности данных и обеспечивает целостность информации. Математическая основа реляционной модели обеспечила строгость и предсказуемость операций с данными, что сделало ее доминирующей в мире информационных технологий на несколько десятилетий \cite{date2003}.

\ManualSubsubsection{Нормализация как метод проектирования баз данных}

Процесс нормализации представляет собой систематический подход к проектированию баз данных, направленный на устранение аномалий при операциях вставки, обновления и удаления данных. Нормализация осуществляется через последовательность нормальных форм, каждая из которых предъявляет все более строгие требования к структуре данных.

Первая нормальная форма (1НФ) требует, чтобы все значения в таблице были атомарными, то есть не содержали составных или повторяющихся элементов. Например, вместо хранения нескольких телефонных номеров в одной ячейке через запятую, каждый номер должен занимать отдельную строку или храниться в связанной таблице. Вторая нормальная форма (2НФ) применяется к таблицам с составными первичными ключами и требует, чтобы все неключевые атрибуты полностью зависели от всего ключа, а не только от его части. Это устраняет ситуации, когда изменение одной части ключа требует изменения не связанных с ней данных \cite{elmasri2016}.

Третья нормальная форма (3НФ) решает проблему транзитивных зависимостей, когда неключевые атрибуты зависят друг от друга, а не непосредственно от первичного ключа. Разделение таких атрибутов по разным таблицам устраняет аномалии обновления, когда изменение одного значения требует корректировки нескольких связанных значений. Дальнейшие нормальные формы (нормальная форма Бойса-Кодда, четвертая и пятая нормальные формы) решают более сложные проблемы проектирования, связанные с многозначными зависимостями и сложными связями между атрибутами.

Практическое значение нормализации выходит за рамки теоретического совершенства. Нормализованная база данных требует меньше места для хранения благодаря устранению дублирования данных. Она также упрощает процессы обновления информации, поскольку каждое изменение нужно вносить только в одном месте. Более того, нормализация снижает вероятность возникновения противоречий в данных и делает структуру базы более понятной для разработчиков и администраторов, что упрощает ее сопровождение и развитие \cite{date2003}.

\ManualSubsubsection{Принципы проектирования современных систем хранения данных}

Современные подходы к проектированию баз данных учитывают не только теоретические принципы нормализации, но и практические требования производительности, масштабируемости и доступности. Одним из ключевых компромиссов при проектировании является баланс между степенью нормализации и производительностью запросов. Чрезмерная нормализация может привести к необходимости выполнения большого количества операций соединения таблиц (JOIN), что снижает скорость выполнения запросов. В таких случаях применяется контролируемая денормализация — преднамеренное введение избыточности данных для ускорения операций чтения.

Другим важным принципом является разделение ответственности между уровнем приложения и уровнем базы данных. Современные практики рекомендуют размещать бизнес-логику преимущественно на уровне приложения, оставляя базе данных функции хранения, целостности и базовой обработки данных. Однако сложные правила целостности и транзакционная логика часто эффективнее реализуются на уровне базы данных через хранимые процедуры, триггеры и ограничения \cite{elmasri2016}.

Принцип минимальных привилегий предполагает, что пользователи и приложения должны получать только те права доступа к данным, которые необходимы для выполнения их функций. Этот принцип реализуется через сложные системы управления доступом, включая ролевую модель (RBAC) и модель на основе атрибутов (ABAC). Принцип отказоустойчивости требует, чтобы система продолжала функционировать или быстро восстанавливалась после аппаратных или программных сбоев, что достигается через репликацию данных, резервное копирование и автоматическое переключение при отказе \cite{gray1993}.

\newpage
\refstepcounter{section}
\addcontentsline{toc}{section}{\protect\numberline{\thesection}Отличие реляционных и нереляционных моделей данных}
\begin{center}\textbf{\thesection\ ОТЛИЧИЕ РЕЛЯЦИОННЫХ И НЕРЕЛЯЦИОННЫХ МОДЕЛЕЙ ДАННЫХ}\end{center}
\vspace*{-1.5em}

\ManualSubsection{Основные принципы различия реляционной и нереляционной моделей}

Современная экосистема систем управления данными характеризуется сосуществованием двух основных парадигм: реляционной, основанной на строгих математических принципах и табличной организации данных, и нереляционной, предлагающей более гибкие, специализированные подходы к хранению информации. Понимание различий между этими подходами и областей их эффективного применения является критически важным для проектирования современных информационных систем.

\ManualSubsubsection{Реляционная модель: строгость и универсальность}

Реляционная модель данных, предложенная Эдгаром Коддом в 1970 году, основана на фундаментальных принципах математической логики и теории множеств. Данные представляются в виде набора таблиц (отношений), где каждая таблица соответствует определенной сущности предметной области. Строки таблицы (кортежи) представляют экземпляры сущности, а столбцы (атрибуты) — их свойства. Важнейшим элементом реляционной модели являются ключи: первичный ключ уникально идентифицирует каждую строку в таблице, а внешний ключ устанавливает связи между таблицами, обеспечивая ссылочную целостность данных \cite{codd1970}.

Сильной стороной реляционной модели является ее строгость и предсказуемость. Все операции с данными описываются в терминах реляционной алгебры и реализуются через стандартизированный язык SQL (Structured Query Language). Транзакционная обработка в реляционных базах данных следует принципу ACID (Atomicity, Consistency, Isolation, Durability), гарантируя атомарность операций, согласованность данных, изолированность параллельных транзакций и долговечность результатов. Эти свойства делают реляционные базы данных идеальным выбором для систем, требующих высокой целостности данных, таких как банковские приложения, системы управления запасами и другие корпоративные информационные системы \cite{haerder1983}.

Однако реляционная модель имеет и ограничения. Жесткая схема данных требует тщательного проектирования до начала эксплуатации системы и затрудняет адаптацию к изменяющимся требованиям. Горизонтальное масштабирование (распределение данных между несколькими серверами) в реляционных системах сложно реализуется и часто приводит к снижению производительности из-за необходимости поддержания согласованности данных между узлами \cite{sadalage2012}.

\ManualSubsubsection{Нереляционные модели: разнообразие и специализация}

Термин NoSQL (Not Only SQL) объединяет широкий спектр систем хранения данных, которые отвергают или дополняют принципы реляционной модели. Возникновение этих систем в конце 2000-х годов было связано с потребностями таких компаний, как Google, Amazon и Facebook, которые столкнулись с ограничениями реляционных СУБД при обработке эксабайтов данных в режиме реального времени. В отличие от единообразия реляционной модели, мир NoSQL характеризуется разнообразием подходов, каждый из которых оптимизирован для решения определенного класса задач \cite{sadalage2012}.

Документные базы данных, такие как MongoDB и Couchbase, хранят данные в виде полуструктурированных документов (обычно в форматах JSON или BSON). Каждый документ содержит всю информацию, относящуюся к определенной сущности, что позволяет избежать сложных операций соединения таблиц. Документные базы хорошо подходят для систем управления контентом, каталогов продукции и других приложений, где данные имеют иерархическую структуру и могут изменяться со временем.

Базы данных типа «ключ-значение», включая Redis и Amazon DynamoDB, представляют собой простейшую модель хранения, где данные извлекаются по уникальному ключу. Эти системы обеспечивают максимальную скорость доступа к данным и идеально подходят для кэширования, хранения сессий пользователей, счетчиков и других сценариев, где важна минимальная задержка. Однако они не поддерживают сложные запросы и связи между данными \cite{sadalage2012}.

Колоночные базы данных, такие как Cassandra и HBase, организуют данные по столбцам, а не по строкам. Это позволяет эффективно сжимать данные и быстро выполнять агрегатные запросы, которые затрагивают только подмножество столбцов. Такие системы идеальны для аналитических нагрузок, систем сбора метрик и других приложений, где требуется обработка больших объемов данных с акцентом на операции чтения.

Графовые базы данных, например Neo4j, специализируются на хранении и обработке сложных взаимосвязей между сущностями. Данные представляются в виде узлов (сущностей), ребер (связей) и свойств (атрибутов). Такая организация позволяет эффективно выполнять запросы, связанные с анализом связей, таких как поиск кратчайшего пути между узлами или выявление сообществ в социальных сетях \cite{sadalage2012}.

\ManualSubsubsection{Сравнительный анализ и гибридные подходы}

Ключевое различие между реляционными и нереляционными системами проявляется в их подходе к схеме данных. Реляционные базы требуют определения схемы до начала работы с данными, что обеспечивает строгость и предсказуемость, но ограничивает гибкость. Нереляционные системы часто используют динамические схемы или вообще обходятся без явного определения структуры данных, что позволяет быстро адаптироваться к изменяющимся требованиям, но может привести к проблемам с целостностью данных.

Принципы обеспечения целостности также различаются. Реляционные системы строго следуют принципу ACID, гарантируя полную согласованность данных даже в условиях параллельного доступа. Многие NoSQL системы придерживаются принципа BASE (Basically Available, Soft state, Eventual consistency), который в обмен на более высокую производительность и доступность допускает временную несогласованность данных в разных узлах кластера, с гарантией, что в конечном счете все узлы придут к согласованному состоянию \cite{ongaro2014}.

Важным различием является и подход к масштабированию. Реляционные базы традиционно масштабируются вертикально (добавление ресурсов к существующему серверу), в то время как NoSQL системы проектируются для горизонтального масштабирования (добавление новых серверов в кластер). Теорема CAP (Consistency, Availability, Partition tolerance) формулирует фундаментальное ограничение распределенных систем: при разделении сети невозможно одновременно обеспечить согласованность, доступность и устойчивость к разделению. Реляционные системы обычно выбирают согласованность и доступность, а многие NoSQL системы — доступность и устойчивость к разделению \cite{ongaro2014}.

Современные тенденции показывают движение к конвергенции двух подходов. Многие реляционные СУБД добавляют поддержку JSON и других полуструктурированных форматов данных, а некоторые NoSQL системы начинают поддерживать SQL-подобные языки запросов и транзакции с гарантиями ACID. Гибридные системы, такие как Google Spanner и CockroachDB, сочетают горизонтальную масштабируемость NoSQL с транзакционной целостностью реляционных баз, предлагая новый компромисс в пространстве проектирования систем хранения данных \cite{corbett2012}.

\newpage
\refstepcounter{section}
\addcontentsline{toc}{section}{\protect\numberline{\thesection}Основные функции и возможности систем управления базами данных}
\begin{center}\textbf{\thesection\ ОСНОВНЫЕ ФУНКЦИИ И ВОЗМОЖНОСТИ СИСТЕМ УПРАВЛЕНИЯ БАЗАМИ ДАННЫХ}\end{center}
\vspace*{-1.5em}
\ManualSubsection{Функции и возможности СУБД}

Системы управления базами данных (СУБД) представляют собой сложные программные комплексы, которые обеспечивают интерфейс между физическим хранением данных и прикладными программами. Современные СУБД выполняют широкий спектр функций, выходящих далеко за рамки простого хранения и извлечения информации. Эти функции можно разделить на несколько ключевых категорий, каждая из которых решает определенные задачи в управлении данными.

\ManualSubsubsection{Управление транзакциями и обеспечение целостности данных}

Одной из фундаментальных функций СУБД является управление транзакциями — логическими единицами работы с базой данных, которые должны выполняться полностью или не выполняться вовсе. Эта функция обеспечивает атомарность операций, что критически важно для поддержания целостности данных в системах, где несколько операций должны выполняться как единое целое. Классический пример — перевод денежных средств между банковскими счетами: операция должна либо полностью завершиться (деньги списываются с одного счета и зачисляются на другой), либо не выполниться вообще, чтобы избежать ситуации, когда деньги списаны, но не зачислены \cite{gray1993}.

Для реализации управления транзакциями СУБД используют различные механизмы управления параллелизмом. Пессимистичные блокировки предотвращают конфликты путем блокировки данных на время выполнения транзакции, но могут снижать производительность при высокой конкуренции за ресурсы. Оптимистичные блокировки, такие как многоверсионное управление конкурентным доступом (MVCC), позволяют нескольким транзакциям одновременно работать с данными, проверяя конфликты только на этапе завершения транзакции. MVCC создает отдельные версии данных для каждой транзакции, что позволяет читающим транзакциям получать непротиворечивый снимок данных без блокировок, значительно повышая производительность в средах с преобладанием операций чтения \cite{gray1993}.

Уровни изоляции транзакций определяют степень защиты от различных аномалий параллельного доступа. Стандарт SQL определяет четыре уровня изоляции: Read Uncommitted (допускает чтение незафиксированных данных), Read Committed (гарантирует чтение только зафиксированных данных), Repeatable Read (обеспечивает повторяемость чтений в рамках одной транзакции) и Serializable (полная изоляция, как если бы транзакции выполнялись последовательно). Выбор уровня изоляции представляет собой компромисс между степенью защиты от аномалий и производительностью системы \cite{haerder1983}.

\ManualSubsubsection{Обеспечение безопасности и контроля доступа}

Современные СУБД реализуют сложные системы безопасности, которые защищают данные от несанкционированного доступа и обеспечивают соответствие регуляторным требованиям. Эти системы основаны на модели дискреционного управления доступом, где права доступа к объектам базы данных определяются для каждого пользователя или роли отдельно.

Аутентификация проверяет идентичность пользователей, пытающихся получить доступ к системе. Современные СУБД поддерживают различные методы аутентификации: от простых паролей до интеграции с корпоративными системами каталогов (LDAP, Active Directory) и многофакторной аутентификации. Авторизация определяет, какие операции разрешено выполнять аутентифицированным пользователям. Ролевая модель доступа (RBAC) позволяет назначать права доступа ролям, а затем назначать роли пользователям, что упрощает управление правами в крупных организациях. Модель доступа на основе атрибутов (ABAC) учитывает дополнительные параметры, такие как время суток, местоположение пользователя или чувствительность данных, предоставляя более гибкую систему контроля доступа \cite{gdpr}.

Аудит доступа к данным позволяет отслеживать все операции, выполняемые с базой данных, что важно как для выявления подозрительной активности, так и для соответствия регуляторным требованиям, таким как GDPR, HIPAA или PCI DSS. Современные СУБД предоставляют средства для настройки детализированного аудита, включая регистрацию успешных и неуспешных попыток доступа, изменений схемы базы данных и доступа к конфиденциальным данным. Шифрование данных защищает информацию как при хранении (шифрование на диске), так и при передаче по сети (TLS), предотвращая утечки данных даже в случае компрометации инфраструктуры хранения \cite{gdpr}.

\ManualSubsubsection{Оптимизация производительности и масштабируемость}

Производительность СУБД зависит от множества факторов, включая эффективность выполнения запросов, управление памятью и дисковым вводом-выводом, а также способность системы масштабироваться при увеличении нагрузки. Современные СУБД реализуют сложные механизмы оптимизации, которые автоматически настраивают параметры системы для достижения максимальной производительности при заданных ограничениях ресурсов.

Оптимизатор запросов анализирует SQL-запросы и выбирает наиболее эффективный план их выполнения. Этот процесс включает определение порядка соединения таблиц, выбор методов доступа к данным (через индексы или полное сканирование таблиц) и оценку стоимости различных вариантов выполнения запроса. Статистика по распределению данных, собираемая СУБД, позволяет оптимизатору принимать обоснованные решения. Современные СУБД также поддерживают адаптивную оптимизацию, которая корректирует планы выполнения запросов на основе фактической производительности \cite{elmasri2016}.

Индексы являются ключевым механизмом ускорения доступа к данным. B-деревья обеспечивают эффективный поиск по диапазонам значений и поддерживают операции сортировки. Хэш-индексы предоставляют максимальную скорость точечного поиска, но не поддерживают диапазонные запросы. Специализированные индексы, такие как GiST (Generalized Search Tree) и GIN (Generalized Inverted Index) в PostgreSQL, оптимизированы для работы со сложными типами данных, включая геопространственную информацию, полнотекстовый поиск и массивы \cite{elmasri2016}.

Масштабируемость СУБД может достигаться как вертикально (увеличение ресурсов отдельного сервера), так и горизонтально (добавление новых серверов в кластер). Репликация данных на несколько серверов позволяет распределить нагрузку операций чтения и обеспечить отказоустойчивость. Шардирование (горизонтальное партиционирование) разделяет большие таблицы на части, распределенные по разным серверам, что позволяет обрабатывать объемы данных, превышающие возможности отдельного сервера. Современные распределенные СУБД, такие как Google Spanner и CockroachDB, обеспечивают глобальное распределение данных с сохранением транзакционной целостности, что позволяет создавать высокодоступные приложения, работающие в масштабе всей планеты \cite{corbett2012}.

\ManualSubsubsection{Резервное копирование и восстановление после сбоев}

Надежность СУБД во многом зависит от эффективности механизмов резервного копирования и восстановления данных. Современные СУБД поддерживают различные стратегии резервного копирования, каждая из которых имеет свои преимущества и ограничения. Полное резервное копирование создает копию всей базы данных и обеспечивает простоту восстановления, но требует значительного времени и места для хранения. Инкрементальное резервное копирование сохраняет только изменения, произошедшие с момента последнего резервного копирования, что сокращает время выполнения операции и объем хранимых данных, но усложняет процесс восстановления. Дифференциальное резервное копирование сохраняет все изменения с момента последнего полного резервного копирования, предлагая компромисс между скоростью выполнения и сложностью восстановления \cite{gray1993}.

Журналирование транзакций (Write-Ahead Logging, WAL) является ключевым механизмом обеспечения долговечности данных и восстановления после сбоев. Перед внесением изменений в основные структуры данных СУБД записывает эти изменения в журнал транзакций. В случае сбоя система может использовать журнал для восстановления согласованного состояния базы данных, повторив зафиксированные транзакции и откатив незавершенные. Современные СУБД также поддерживают восстановление на определенный момент времени (Point-in-Time Recovery, PITR), которое позволяет восстановить базу данных в состояние на любой момент в пределах периода, охватываемого архивными копиями журналов транзакций \cite{gray1993}.

Метрики RPO (Recovery Point Objective) и RTO (Recovery Time Objective) определяют требования к системе восстановления. RPO указывает максимально допустимый объем данных, который может быть потерян при аварии, что определяет необходимую частоту резервного копирования. RTO устанавливает максимально допустимое время простоя системы после сбоя, что влияет на архитектуру отказоустойчивости и процедуры восстановления. Для критических систем, таких как банковские приложения, RPO и RTO могут стремиться к нулю, что требует реализации сложных и дорогостоящих решений с синхронной репликацией данных и мгновенным переключением при отказе \cite{gray1993}.

\newpage
\refstepcounter{section}
\addcontentsline{toc}{section}{\protect\numberline{\thesection}Язык SQL как инструмент взаимодействия с базами данных}
\begin{center}\textbf{\thesection\ ЯЗЫК SQL КАК ИНСТРУМЕНТ ВЗАИМОДЕЙСТВИЯ С БАЗАМИ ДАННЫХ}\end{center}
\vspace*{-1.5em}

\ManualSubsection{Роль SQL в работе с базами данных}

Язык SQL (Structured Query Language) представляет собой стандартизированный язык программирования, предназначенный для управления данными в реляционных базах данных. Разработанный в 1970-х годах в IBM, SQL прошел значительную эволюцию и стал де-факто стандартом для работы с базами данных, поддерживаемым всеми основными реляционными СУБД. Декларативная природа SQL позволяет пользователям описывать, какие данные необходимо получить или изменить, оставляя за СУБД задачу определения наиболее эффективного способа выполнения операции \cite{date2003}.

\ManualSubsubsection{Архитектура и компоненты языка SQL}

SQL состоит из нескольких подъязыков, каждый из которых отвечает за определенный аспект работы с базой данных. Язык определения данных (DDL — Data Definition Language) включает операторы для создания, изменения и удаления структурных элементов базы данных. CREATE TABLE определяет структуру новой таблицы, включая имена и типы столбцов, ограничения целостности и связи с другими таблицами. ALTER TABLE позволяет изменять существующую структуру таблицы, а DROP TABLE — удалять таблицы из базы данных. Аналогичные операторы существуют для других объектов базы данных: индексов (CREATE INDEX, DROP INDEX), представлений (CREATE VIEW, DROP VIEW) и хранимых процедур (CREATE PROCEDURE, DROP PROCEDURE) \cite{elmasri2016}.

Язык манипулирования данными (DML — Data Manipulation Language) содержит операторы для работы с содержимым таблиц. SELECT извлекает данные из одной или нескольких таблиц согласно заданным критериям. INSERT добавляет новые строки в таблицу, UPDATE изменяет существующие данные, а DELETE удаляет строки из таблицы. Язык управления данными (DCL — Data Control Language) включает операторы для управления доступом к данным: GRANT предоставляет права доступа пользователям или ролям, а REVOKE отзывает ранее предоставленные права. Язык управления транзакциями (TCL — Transaction Control Language) содержит операторы для управления транзакциями: BEGIN TRANSACTION отмечает начало транзакции, COMMIT фиксирует изменения, а ROLLBACK отменяет все изменения, сделанные в текущей транзакции \cite{elmasri2016}.

\ManualSubsubsection{Расширенные возможности SQL для сложных запросов}

Одной из мощнейших возможностей SQL является выполнение сложных запросов, объединяющих данные из нескольких таблиц. Оператор JOIN позволяет связать таблицы по ключевым полям, формируя результирующий набор, который включает информацию из всех участвующих таблиц. Внутреннее соединение (INNER JOIN) возвращает только те записи, для которых найдены соответствия во всех таблицах. Левое внешнее соединение (LEFT JOIN) возвращает все записи из левой таблицы и соответствующие записи из правой таблицы, или NULL, если соответствие не найдено. Правое внешнее соединение (RIGHT JOIN) работает аналогично, но возвращает все записи из правой таблицы. Полное внешнее соединение (FULL JOIN) возвращает все записи из обеих таблиц, заполняя отсутствующие значения NULL \cite{date2003}.

Агрегатные функции позволяют вычислять сводные показатели по группам данных. SUM вычисляет сумму значений, AVG — среднее значение, COUNT — количество записей, MIN и MAX — минимальное и максимальное значения. В сочетании с предложением GROUP BY эти функции вычисляют агрегированные значения для каждой группы записей. Предложение HAVING позволяет фильтровать группы по результатам агрегации. Оконные функции (window functions) расширяют возможности агрегации, позволяя вычислять значения для каждой строки результирующего набора на основе подмножества строк (окна) без свертки строк в группы. Это особенно полезно для вычисления скользящих средних, ранжирования и анализа временных рядов \cite{date2003}.

Рекурсивные запросы (WITH RECURSIVE) позволяют работать с иерархическими структурами данных, такими как организационные диаграммы, деревья категорий или графы. Общие табличные выражения (CTE — Common Table Expressions) позволяют определять временные результирующие наборы, которые можно использовать в основном запросе, улучшая читаемость сложных запросов и поддерживая рекурсию. Полнотекстовый поиск предоставляет возможности для эффективного поиска по текстовым полям с учетом морфологии языка, релевантности и других параметров \cite{postgresqldocs}.

\ManualSubsubsection{Диалекты SQL и интеграция с современными технологиями}

Несмотря на существование стандартов SQL (ANSI SQL, ISO/IEC 9075), производители СУБД развивают собственные диалекты языка, добавляя уникальные возможности и оптимизации. PostgreSQL предлагает расширенные типы данных для работы с географической информацией (PostGIS), массивами, диапазонами значений и составными типами. Его диалект SQL включает мощные возможности для работы с JSON, оконные функции с расширенным синтаксисом и поддержку рекурсивных запросов. Oracle SQL отличается развитыми аналитическими функциями для бизнес-анализа, расширенными возможностями для работы с иерархическими данными и интеграцией с языком PL/SQL для процедурного программирования. Microsoft SQL Server предоставляет собственную реализацию SQL с интеграцией в платформу .NET, поддержкой выполнения кода на CLR внутри базы данных и развитыми средствами для бизнес-аналитики. MySQL, будучи одной из самых популярных СУБД для веб-приложений, оптимизирована для сценариев с высокой интенсивностью чтения и предлагает простой, но эффективный диалект SQL \cite{postgresqldocs}.

Интеграция SQL с современными технологиями значительно расширяет область его применения. В экосистеме больших данных инструменты вроде Apache Hive и Spark SQL позволяют использовать знакомый синтаксис SQL для работы с распределенными хранилищами данных, такими как Hadoop HDFS. Встраиваемые СУБД, такие как SQLite, предоставляют возможности SQL внутри мобильных и десктопных приложений без необходимости развертывания отдельного сервера базы данных. Поддержка стандарта SQL в облачных сервисах (Amazon RDS, Google Cloud SQL, Azure SQL Database) делает язык универсальным инструментом для работы с данными независимо от среды выполнения. Граничные вычисления и IoT устройства все чаще включают легковесные СУБД с поддержкой SQL для локальной обработки данных перед их передачей в централизованные хранилища \cite{tpc}.

\newpage
\refstepcounter{section}
\addcontentsline{toc}{section}{\protect\numberline{\thesection}Концепции целостности и надежности в системах хранения данных}
\begin{center}\textbf{\thesection\ КОНЦЕПЦИИ ЦЕЛОСТНОСТИ И НАДЕЖНОСТИ В СИСТЕМАХ ХРАНЕНИЯ ДАННЫХ}\end{center}
\vspace*{-1.5em}

\ManualSubsection{Целостность и надежность данных в системах хранения}

Целостность и надежность данных являются фундаментальными требованиями к современным системам хранения информации. Эти концепции охватывают широкий спектр аспектов — от гарантий корректности отдельных операций до обеспечения непрерывной работы системы в условиях аппаратных и программных сбоев. Реализация этих концепций требует комплексного подхода, сочетающего технические решения, организационные процедуры и соответствие регуляторным требованиям.

\ManualSubsubsection{Уровни целостности данных в информационных системах}

Целостность данных в системах хранения информации обеспечивается на нескольких взаимосвязанных уровнях. Сущностная целостность гарантирует уникальность каждой записи в таблице через механизмы первичных ключей и ограничений уникальности. Первичный ключ однозначно идентифицирует каждую строку таблицы, предотвращая дублирование записей. Ограничение UNIQUE обеспечивает уникальность значений в столбце или группе столбцов, даже если они не являются первичным ключом. Эти механизмы реализуются как на уровне декларативных ограничений в определении таблицы, так и через индексы, которые обеспечивают эффективную проверку уникальности при операциях вставки и обновления данных \cite{haerder1983}.

Ссылочная целостность поддерживает корректность связей между таблицами через внешние ключи. Внешний ключ в дочерней таблице ссылается на первичный ключ в родительской таблице, предотвращая появление «висячих» ссылок на несуществующие записи. Каскадные операции определяют поведение системы при удалении или обновлении связанных записей. ON DELETE CASCADE автоматически удаляет дочерние записи при удалении родительской, ON DELETE SET NULL устанавливает значение внешнего ключа в NULL, а ON DELETE RESTRICT запрещает удаление родительской записи, если существуют связанные дочерние записи. Аналогичные варианты поведения определяются для операции обновления (ON UPDATE) \cite{elmasri2016}.

Семантическая целостность обеспечивает соответствие данных бизнес-правилам и ограничениям предметной области. Ограничение CHECK позволяет задавать произвольные условия, которым должны удовлетворять значения в столбце или строке. Например, можно определить, что возраст сотрудника должен быть больше 18 лет, или что дата окончания проекта не может быть раньше даты его начала. Ограничение NOT NULL гарантирует, что столбец не содержит пустых значений. Триггеры — автоматически выполняемые процедуры при определенных событиях (вставка, обновление, удаление данных) — позволяют реализовать сложную бизнес-логику проверки и преобразования данных \cite{haerder1983}.

\ManualSubsubsection{Теорема CAP и модели согласованности в распределенных системах}

Теорема CAP, сформулированная Эриком Брюэром в 2000 году, устанавливает фундаментальное ограничение для распределенных систем хранения данных. Согласно этой теореме, при разделении сети (Partition tolerance) невозможно одновременно обеспечить согласованность данных (Consistency) и доступность системы (Availability). Согласованность означает, что все узлы системы видят одни и те же данные в один момент времени. Доступность гарантирует, что каждый запрос к системе получает ответ (успешный или неуспешный). Устойчивость к разделению означает, что система продолжает функционировать при разрыве связи между узлами \cite{ongaro2014}.

На практике системы выбирают компромисс между этими тремя свойствами. Реляционные базы данных, развернутые на одном сервере или в кластере с синхронной репликацией, обычно выбирают согласованность и доступность (CA), но оказываются уязвимыми к разделению сети. Многие NoSQL системы, такие как Cassandra и DynamoDB, предпочитают доступность и устойчивость к разделению (AP), обеспечивая согласованность в конечном счете (eventual consistency). Специализированные распределенные системы, такие как Google Spanner и CockroachDB, пытаются преодолеть ограничения теоремы CAP через использование точных атомарных часов и специальных алгоритмов консенсуса, предлагая сильную согласованность при сохранении высокой доступности и устойчивости к разделению \cite{corbett2012}.

Модели согласованности определяют, какие гарантии предоставляет система относительно актуальности данных. Сильная согласованность (strong consistency) гарантирует, что все узлы видят одни и те же данные в один момент времени. Согласованность в конечном счете (eventual consistency) обещает, что если в систему не вносятся новые изменения, то через некоторое время все узлы придут к согласованному состоянию. Согласованность на чтение (read-your-writes consistency) гарантирует, что пользователь всегда видит свои собственные обновления. Согласованность по сессиям (session consistency) обеспечивает согласованность в рамках одной сессии пользователя. Выбор модели согласованности зависит от требований приложения: финансовые системы требуют сильной согласованности, тогда как социальные сети могут довольствоваться согласованностью в конечном счете \cite{ongaro2014}.

\ManualSubsubsection{Распределенные транзакции и протоколы консенсуса}

Реализация транзакций в распределенных системах представляет собой сложную задачу, требующую координации действий между несколькими независимыми узлами. Двухфазный коммит (2PC — Two-Phase Commit) является классическим протоколом для реализации распределенных транзакций. В первой фазе координатор запрашивает все участников подготовиться к выполнению транзакции. Если все участники готовы, координатор во второй фазе отправляет команду на выполнение транзакции. Если хотя бы один участник не готов или происходит сбой, координатор отменяет транзакцию. Несмотря на свою концептуальную простоту, 2PC имеет серьезные недостатки: блокировка ресурсов на время выполнения протокола, координатор как единая точка отказа и сложность восстановления после сбоев \cite{gray1993}.

Альтернативные подходы к реализации распределенных транзакций включают паттерн SAGA и протокол TCC (Try-Confirm-Cancel). SAGA разбивает распределенную транзакцию на последовательность локальных транзакций, каждая из которых имеет компенсирующую транзакцию для отката изменений в случае сбоя. TCC использует трехфазный подход: сначала участники резервируют ресурсы (Try), затем подтверждают их использование (Confirm) или отменяют резервирование (Cancel). Эти подходы лучше подходят для микросервисных архитектур, где различные сервисы управляют своими собственными базами данных \cite{gray1993}.

Протоколы консенсуса, такие как Paxos и Raft, решают задачу достижения согласия между распределенными узлами в условиях возможных сбоев. Paxos, предложенный Лесли Лэмпортом в 1990 году, стал теоретической основой для многих распределенных систем. Raft, разработанный в 2014 году, предлагает более понятную альтернативу Paxos с акцентом на понятность и реализуемость. Эти протоколы лежат в основе распределенных СУБД, систем управления конфигурацией (ZooKeeper, etcd) и блокчейн-систем, обеспечивая согласованность данных между узлами даже при сбоях отдельных участников \cite{lamport1998, ongaro2014}.

\newpage
\refstepcounter{section}
\addcontentsline{toc}{section}{\protect\numberline{\thesection}Облачные подходы и онлайн-инструменты для работы с базами данных}
\begin{center}\textbf{\thesection\ ОБЛАЧНЫЕ ПОДХОДЫ И ОНЛАЙН-ИНСТРУМЕНТЫ ДЛЯ РАБОТЫ С БАЗАМИ ДАННЫХ}\end{center}
\vspace*{-1.5em}

\ManualSubsection{Облачные БД и онлайн-инструменты для работы с данными}

Переход к облачным вычислениям кардинально изменил подходы к развертыванию, управлению и использованию баз данных. Облачные провайдеры предлагают широкий спектр услуг в области управления данными, от простых управляемых инстансов традиционных СУБД до специализированных сервисов для конкретных сценариев использования. Эти услуги позволяют организациям сосредоточиться на разработке приложений и анализе данных, передавая рутинные операции администрирования, обновления, резервного копирования и масштабирования облачному провайдеру.

\ManualSubsubsection{Модель «база данных как сервис» (DBaaS) и ее преимущества}

Модель «база данных как сервис» (DBaaS — Database as a Service) представляет собой облачную услугу, которая предоставляет доступ к функциональности базы данных без необходимости управления инфраструктурой, установки программного обеспечения или выполнения рутинных задач администрирования. Ведущие поставщики облачных услуг, включая Amazon Web Services (AWS), Google Cloud Platform (GCP) и Microsoft Azure, предлагают широкий спектр DBaaS решений как для реляционных, так и для нереляционных баз данных \cite{cockroach}.

Основные преимущества DBaaS включают экономию затрат, масштабируемость и снижение операционной нагрузки. Организации могут избежать значительных капитальных затрат на приобретение оборудования и лицензий на программное обеспечение, перейдя на модель оплаты по мере использования. Автоматическое масштабирование позволяет динамически увеличивать или уменьшать вычислительные ресурсы в зависимости от текущей нагрузки, что обеспечивает оптимальное использование ресурсов и снижение затрат. Провайдеры DBaaS берут на себя рутинные задачи администрирования, такие как установка обновлений безопасности, настройка резервного копирования, мониторинг производительности и восстановление после сбоев, что позволяет ИТ-командам организаций сосредоточиться на стратегических задачах \cite{clickhousedocs}.

DBaaS также обеспечивает более высокий уровень доступности и надежности по сравнению с традиционными локальными развертываниями. Облачные провайдеры размещают данные в географически распределенных центрах обработки данных с избыточной инфраструктурой, что обеспечивает защиту от локальных сбоев и катастроф. Автоматическое резервное копирование и восстановление на определенный момент времени (PITR) минимизируют риск потери данных. Встроенные механизмы репликации и автоматического переключения при отказе (failover) обеспечивают высокую доступность сервиса даже в случае сбоев оборудования или программного обеспечения \cite{cockroach}.

\ManualSubsubsection{Облачные сервисы для различных типов баз данных}

Облачные провайдеры предлагают управляемые сервисы для всех основных типов баз данных. Для реляционных СУБД Amazon RDS (Relational Database Service) поддерживает MySQL, PostgreSQL, Oracle, SQL Server, и MariaDB. Google Cloud SQL предоставляет управляемые сервисы для MySQL, PostgreSQL и SQL Server. Azure SQL Database предлагает управляемый сервис на основе Microsoft SQL Server с дополнительными возможностями, такими как встроенный искусственный интеллект для оптимизации производительности. Российские провайдеры, такие как Yandex Cloud, также предлагают управляемые сервисы для популярных открытых СУБД \cite{clickhousedocs}.

Для нереляционных баз данных облачные провайдеры предлагают специализированные сервисы. Документные базы представлены Amazon DocumentDB (совместимый с MongoDB API) и Azure Cosmos DB (мультимодельная БД с поддержкой нескольких API). Базы данных «ключ-значение» реализованы через Amazon ElastiCache для Redis и Memcached, а также Azure Cache for Redis. Колоночные базы доступны через Amazon Keyspaces (совместимый с Apache Cassandra), а графовые — через Amazon Neptune и графовый API Azure Cosmos DB. Google Cloud предлагает Firestore — гибкую, масштабируемую документную базу данных для мобильных, веб- и серверных разработок \cite{kafka}.

Специализированные сервисы для аналитических нагрузок включают Amazon Redshift (управляемый хранилище данных на основе колоночной СУБД), Google BigQuery (бессерверное хранилище данных с поддержкой SQL-запросов к петабайтам данных) и Azure Synapse Analytics (интегрированная аналитическая служба, объединяющая хранилище данных большого объема и аналитику больших данных). Эти сервисы оптимизированы для выполнения сложных аналитических запросов к большим объемам данных и часто интегрируются с экосистемами машинного обучения и бизнес-аналитики соответствующих облачных платформ \cite{clickhousedocs}.

\ManualSubsubsection{Онлайн-инструменты для разработки и администрирования баз данных}

Современные облачные платформы предоставляют богатый набор онлайн-инструментов для администрирования, мониторинга и разработки баз данных. Веб-консоли облачных провайдеров позволяют управлять инстансами баз данных, настраивать параметры, просматривать метрики производительности (использование CPU, памяти, дискового ввода-вывода) и настраивать оповещения о проблемах. Эти консоли также предоставляют доступ к журналам ошибок и диагностической информации, что упрощает устранение неполадок \cite{prometheus}.

Инструменты мониторинга, такие как Amazon CloudWatch, Google Cloud Monitoring и Azure Monitor, предоставляют детализированную информацию о работе баз данных, включая метрики производительности, использование ресурсов и активность подключений. Эти инструменты позволяют настраивать дашборды для визуализации ключевых показателей, создавать оповещения о превышении пороговых значений и анализировать тенденции изменения нагрузки. Интеграция с системами управления инцидентами, такими как PagerDuty или OpsGenie, позволяет автоматизировать реагирование на проблемы с базами данных \cite{prometheus}.

Для разработчиков доступны различные онлайн-среды и инструменты. Классические веб-интерфейсы, такие как phpMyAdmin для MySQL и pgAdmin для PostgreSQL, могут быть развернуты в облаке для удобного доступа к управлению базами данных. Современные кросс-платформенные клиенты, включая DataGrip от JetBrains, Beekeeper Studio и DBeaver, поддерживают подключение к облачным базам данных и предлагают продвинутые возможности для разработки запросов, отладки и визуализации схем данных. Интегрированные среды разработки, такие как Visual Studio Code с расширениями для работы с базами данных, позволяют эффективно работать с облачными БД непосредственно из редактора кода \cite{debezium}.

Средства автоматизации и управления инфраструктурой как код (IaC) позволяют программировать развертывание и настройку баз данных. Terraform и AWS CloudFormation позволяют описывать инфраструктуру баз данных в декларативных конфигурационных файлах, что обеспечивает воспроизводимость развертываний и контроль версий конфигурации. Инструменты непрерывной интеграции и непрерывного развертывания (CI/CD), такие как Jenkins, GitLab CI и GitHub Actions, могут автоматизировать процессы миграции схемы базы данных, тестирования производительности и развертывания изменений в различных средах \cite{debezium}.

\ManualSubsubsection{Бессерверные базы данных и будущее облачных технологий управления данными}

Бессерверные (serverless) базы данных представляют собой следующую эволюционную ступень в развитии облачных услуг управления данными. В отличие от традиционных управляемых сервисов, где пользователи оплачивают выделенные вычислительные ресурсы, бессерверные базы взимают плату только за фактически выполненные операции и потребленное хранилище. Это позволяет организациям значительно снизить затраты на разработку и эксплуатацию приложений с переменной или непредсказуемой нагрузкой \cite{cockroach}.

Amazon Aurora Serverless, Azure SQL Database serverless и Google Cloud Spanner autoscaler являются примерами бессерверных реляционных баз данных. Эти сервисы автоматически масштабируют вычислительные ресурсы в зависимости от текущей нагрузки, обеспечивая высокую производительность при пиковых нагрузках и минимальные затраты в периоды низкой активности. Бессерверные NoSQL базы, такие как Amazon DynamoDB on-demand и Azure Cosmos DB serverless, предлагают аналогичные возможности для нереляционных моделей данных \cite{cockroach}.

Будущее облачных технологий управления данными связано с дальнейшей интеграцией искусственного интеллекта и машинного обучения. Облачные провайдеры уже внедряют AI-оптимизаторы, которые автоматически анализируют шаблоны запросов и настраивают индексы, параметры и планы выполнения для достижения максимальной производительности. Прогностическое масштабирование позволяет системам заранее увеличивать ресурсы перед ожидаемыми пиками нагрузки. Автоматическое исправление проблем (self-healing) позволяет системам самостоятельно выявлять и устранять сбои без вмешательства человека \cite{cockroach}.

Конвергентные базы данных, сочетающие возможности реляционных и нереляционных моделей в единой платформе, позволяют разработчикам выбирать наиболее подходящую модель данных для каждой задачи без необходимости использования нескольких систем. Глобально распределенные базы данных с гарантированной согласованностью и низкой задержкой обеспечивают поддержку приложений, пользователи которых распределены по всему миру. Интеграция блокчейн-технологий с традиционными системами управления данными открывает новые возможности для создания приложений, требующих неизменяемого аудита операций и децентрализованного хранения данных \cite{corbett2012}.

\newpage
\section*{}
\vspace*{-2.5em}
\begin{center}\textbf{ЗАКЛЮЧЕНИЕ}\end{center}
\phantomsection
\addcontentsline{toc}{section}{Заключение}

Технологии баз данных представляют собой сложную и динамично развивающуюся экосистему, которая служит фундаментом для современных информационных систем. Проведенное исследование показало, что эффективная работа с данными требует не только технических навыков использования конкретных инструментов, но и глубокого понимания архитектурных принципов, компромиссов и стратегических подходов к организации информации.

Фундаментом проектирования любых систем хранения данных являются принципы структурирования информации, включая процесс нормализации, который обеспечивает устранение избыточности, поддержание целостности и создание логичной, понятной структуры. Эти принципы остаются актуальными независимо от выбранной модели данных, формируя основу для эффективного управления информацией на протяжении всего ее жизненного цикла.

Современный ландшафт технологий баз данных характеризуется сосуществованием реляционных и нереляционных подходов, каждый из которых оптимизирован для определенного класса задач. Реляционные системы, с их строгой математической основой, транзакционными гарантиями ACID и мощным языком SQL, остаются незаменимыми для приложений, требующих высокой целостности данных и сложных запросов. Нереляционные системы предлагают альтернативные парадигмы — горизонтальную масштабируемость, гибкость схемы данных и специализацию для конкретных сценариев использования, что делает их идеальным выбором для современных веб-приложений, аналитических систем и работы с большими данными. Ключевым результатом исследования является понимание, что выбор между этими подходами должен основываться на тщательном анализе требований конкретного приложения, а не следовании модным тенденциям.

Язык SQL подтвердил свой статус универсального стандарта для работы со структурированными данными, постоянно развиваясь и адаптируясь к современным требованиям. Его декларативная природа, богатые возможности для сложных запросов и аналитических операций, а также интеграция с разнообразными технологиями делают SQL критически важным инструментом в арсенале любого специалиста по данным. При этом развитие диалектов SQL в различных СУБД и их интеграция с экосистемами больших данных и облачными платформами демонстрируют жизнеспособность и адаптивность этого языка.

Концепции целостности и надежности данных приобретают особую важность в условиях роста объемов информации и распределенности систем. Теорема CAP формулирует фундаментальные ограничения распределенных систем, определяя необходимость компромиссов между согласованностью, доступностью и устойчивостью к разделению. Современные подходы к обеспечению надежности, включая репликацию, шардирование, резервное копирование и стратегии восстановления, позволяют создавать системы, соответствующие строгим требованиям современных приложений. Особое значение имеют метрики RPO и RTO, которые количественно определяют требования к системе восстановления и влияют на архитектурные решения.

Переход к облачным моделям работы с базами данных представляет собой одну из наиболее значимых тенденций в области управления информацией. Модель «база данных как сервис» (DBaaS) кардинально меняет экономику и операционную модель использования баз данных, предлагая экономическую эффективность, автоматическое масштабирование и снижение операционной нагрузки. Бессерверные базы данных и интеграция искусственного интеллекта для оптимизации производительности указывают на направления дальнейшего развития облачных технологий управления данными.

Таким образом, современные технологии баз данных образуют многоуровневую экосистему, где успех зависит от комплексного понимания взаимосвязей между различными компонентами: принципами организации данных, архитектурными моделями, языками запросов, механизмами обеспечения целостности и надежности, а также моделями развертывания и эксплуатации. Освоение этой экосистемы формирует ключевую компетенцию для специалистов в области информационных технологий, позволяя не только эффективно решать текущие задачи, но и создавать масштабируемые, надежные и адаптивные системы, способные отвечать вызовам цифровой трансформации.

\newpage
\begin{center}\textbf{СПИСОК СОКРАЩЕНИЙ}\end{center}
\phantomsection
\addcontentsline{toc}{section}{Список сокращений}

\vspace{1em}

{\small
\noindent
\begin{tabular}{|p{0.15\textwidth}|p{0.80\textwidth}|}
\hline
\textbf{Сокращение} & \multicolumn{1}{c|}{\textbf{Расшифровка}} \\
\noalign{\hrule height 1.2pt}
\hline
БД & База данных \\
\hline
СУБД & Система управления базами данных \\
\hline
SQL & Structured Query Language (язык структурированных запросов) \\
\hline
NoSQL & Not Only SQL (не только SQL) \\
\hline
ACID & Atomicity, Consistency, Isolation, Durability (атомарность, согласованность, изолированность, долговечность) \\
\hline
BASE & Basically Available, Soft state, Eventual consistency (базовая доступность, мягкое состояние, согласованность в конечном счете) \\
\hline
CAP & Consistency, Availability, Partition tolerance (теорема CAP: согласованность, доступность, устойчивость к разделению) \\
\hline
DBaaS & Database as a Service (база данных как сервис) \\
\hline
RPO & Recovery Point Objective (целевая точка восстановления) \\
\hline
RTO & Recovery Time Objective (целевое время восстановления) \\
\hline
API & Application Programming Interface (интерфейс прикладного программирования) \\
\hline
ETL & Extract, Transform, Load (извлечение, преобразование, загрузка) \\
\hline
ELT & Extract, Load, Transform (извлечение, загрузка, преобразование) \\
\hline
CDC & Change Data Capture (отслеживание изменений данных) \\
\hline
OLTP & Online Transaction Processing (оперативная обработка транзакций) \\
\hline
OLAP & Online Analytical Processing (оперативная аналитическая обработка) \\
\hline
JSON & JavaScript Object Notation (обозначение объектов JavaScript) \\
\hline
XML & eXtensible Markup Language (расширяемый язык разметки) \\
\hline
CRUD & Create, Read, Update, Delete (создание, чтение, обновление, удаление) \\
\hline
IDE & Integrated Development Environment (интегрированная среда разработки) \\
\hline
MVCC & Multi-Version Concurrency Control (многовариантное управление параллелизмом) \\
\hline
WAL & Write-Ahead Logging (журнал предзаписи) \\
\hline
DDL & Data Definition Language (язык определения данных) \\
\hline
DML & Data Manipulation Language (язык манипулирования данными) \\
\hline
DCL & Data Control Language (язык управления данными) \\
\hline
TCL & Transaction Control Language (язык управления транзакциями) \\
\hline
BI & Business Intelligence (бизнес-аналитика) \\
\hline
DW & Data Warehouse (хранилище данных) \\
\hline
PK & Primary Key (первичный ключ) \\
\hline
FK & Foreign Key (внешний ключ) \\
\hline
B-Tree & Balanced Tree (сбалансированное дерево) \\
\hline
BIFF & Binary Interchange File Format (двоичный формат обмена файлами) \\
\hline
\end{tabular}
}

\clearpage
\nocite{*}
\printbibliography[heading=bibliography]

\end{document}
